%===============================================================================
% Classic resume template (FreeMarker + LaTeX)
%-------------------------------------------------------------------------------
% Same layout as resume-classic.tex (the filled example), driven by data.
%
% Written to work in BOTH:
%   * real LaTeX engines (pdflatex / xelatex / Overleaf), and
%   * latex.js based web viewers, which only support a small subset of LaTeX
%     (no \newcommand, no \if, no tabular, no \hfill, no \rule, no enumitem
%     options, no titlesec / fancyhdr).
% The raw file compiles too: it then shows the placeholders as text.
% (Never write a bracket-equals or bracket-hash pair in a comment: FreeMarker
% would parse it as a placeholder or directive.)
%
% It is an Apache FreeMarker (2.3.28+) template. The backend must use the
% square-bracket syntax and UTF-8:
%
%   cfg.setTagSyntax(Configuration.SQUARE_BRACKET_TAG_SYNTAX);
%   cfg.setInterpolationSyntax(Configuration.SQUARE_BRACKET_INTERPOLATION_SYNTAX);
%   cfg.setDefaultEncoding("UTF-8");
%   cfg.getTemplate("classic.tex").process(resumeDocument, writer);
%
% Data model root: com.learning.resumemaker.document.ResumeDocument
%   profile (alias p) : fullName, jobTitle, email, phone, location, linkedin, summary
%   education[]       : university, location, degree, grade, startYear, endYear
%   experience[]      : jobTitle, company, startDate, endDate, current, achievements
%   projects[]        : name, technologies, description
%   certifications[]  : name, issuer
%   skills (alias s)  : languages[], frameworks[], tools[]
%   spokenLanguages[] : language, proficiency
%
% Rules followed by this file:
%   * Every FreeMarker directive sits on its own line that starts with "%", so
%     LaTeX and latex.js read it as a comment. A directive never spans lines.
%   * Every placeholder sits inside its own { } group. latex.js rejects a "["
%     directly inside a macro argument, and real LaTeX would read "\item [" or
%     "\\[" as an optional argument.
%   * Bold / italic around data uses {\bfseries ...} / {\itshape ...}, not
%     \textbf / \textit: latex.js crashes on \textbf{} with an empty value.
%   * Text values go through tex() (LaTeX escaping), link targets through url().
%   * Empty values are skipped by FreeMarker, so no LaTeX logic is needed.
%   * experience.achievements and projects.description are split on new lines;
%     each non-empty line becomes one bullet.
% Page geometry: letter paper, 0.5in margins -> text width 7.5in (542pt).
%===============================================================================
%
% [#-- Escape LaTeX special characters and collapse line breaks to spaces. --]
% [#function tex value=""]
% [#local v = value?string?replace("\\", "\x0001")?replace("{", r"\{")?replace("}", r"\}")]
% [#local v = v?replace("&", r"\&")?replace("%", r"\%")?replace("$", r"\$")?replace("#", r"\#")?replace("_", r"\_")]
% [#local v = v?replace("~", r"\textasciitilde{}")?replace("^", r"\textasciicircum{}")]
% [#local v = v?replace("<", r"$<$")?replace(">", r"$>$")?replace("|", r"$|$")]
% [#return v?replace("\x0001", r"\textbackslash{}")?replace(r"\s+", " ", "r")?trim]
% [/#function]
%
% [#-- Make a value safe to use as an \href target. --]
% [#function url value=""]
% [#return value?string?trim?replace(r"[\\{}]", "", "r")?replace(r"%(?![0-9A-Fa-f]{2})", "%25", "r")?replace("%", r"\%")?replace(r"\s+", " ", "r")?replace(" ", r"\%20")]
% [/#function]
%
% [#-- Full LinkedIn link, adding https:// when missing. --]
% [#function linkedinUrl value=""]
% [#local v = value?string?trim]
% [#if v?matches(r"(?i)^https?://.*")][#return v][/#if]
% [#return "https://" + v]
% [/#function]
%
% [#-- LinkedIn text shown on the page, without protocol / www / trailing slash. --]
% [#function linkedinText value=""]
% [#return value?string?trim?replace(r"(?i)^https?://(www\.)?", "", "r")?remove_ending("/")]
% [/#function]
%
% [#-- Escape each non-empty part and join them, e.g. ["Jan 2024", "Present"] -> "Jan 2024 -- Present". --]
% [#function joinTex parts sep]
% [#local out = []]
% [#list parts as part][#if part?has_content][#local out = out + [tex(part)]][/#if][/#list]
% [#return out?join(sep)]
% [/#function]
%
% [#-- End date of a job: "Present" for the current one. --]
% [#function endLabel job]
% [#if job.current!false][#return "Present"][/#if]
% [#return job.endDate!""]
% [/#function]
%
% [#-- Split a multi-line string into bullet lines (drops blanks and leading -, *, bullet signs). --]
% [#function bullets text=""]
% [#local items = []]
% [#list text?string?split(r"\r?\n", "r") as raw]
% [#local line = raw?trim?replace(r"^[-*•·]+\s*", "", "r")]
% [#if line?has_content][#local items = items + [line]][/#if]
% [/#list]
% [#return items]
% [/#function]
%
% [#-- "English (Fluent), Hindi (Native)", already escaped. --]
% [#function spokenList items=[]]
% [#local parts = []]
% [#list items as sl]
% [#if (sl.language)?has_content]
% [#if (sl.proficiency)?has_content]
% [#local parts = parts + [tex(sl.language) + " (" + tex(sl.proficiency) + ")"]]
% [#else]
% [#local parts = parts + [tex(sl.language)]]
% [/#if]
% [/#if]
% [/#list]
% [#return parts?join(", ")]
% [/#function]
%
% [#assign p = profile!{}]
% [#assign s = skills!{}]
%
% [#-- Skill rows: label + escaped value, only the non-empty ones. --]
% [#assign skillRows = []]
% [#if (s.languages)?has_content][#assign skillRows = skillRows + [{"label": "Programming Languages", "value": joinTex(s.languages, ", ")}]][/#if]
% [#if (s.frameworks)?has_content][#assign skillRows = skillRows + [{"label": "Frameworks", "value": joinTex(s.frameworks, ", ")}]][/#if]
% [#if (s.tools)?has_content][#assign skillRows = skillRows + [{"label": "Tools/Technologies", "value": joinTex(s.tools, ", ")}]][/#if]
% [#assign spoken = spokenList(spokenLanguages![])]
% [#if spoken?has_content][#assign skillRows = skillRows + [{"label": "Spoken Languages", "value": spoken}]][/#if]
%-------------------------------------------------------------------------------
\documentclass[11pt]{article}

\usepackage[letterpaper,margin=0.5in]{geometry}
\usepackage[hidelinks]{hyperref}

\pagestyle{empty}
\setlength{\textwidth}{7.5in}
\setlength{\unitlength}{1pt}

\begin{document}

%----------HEADING----------
\begin{center}
% [#if (p.fullName)?has_content]
{\Huge\bfseries {[=tex(p.fullName)]}}\\[6pt]
% [/#if]
% [#if (p.jobTitle)?has_content]
{\large {[=tex(p.jobTitle)]}}\\[4pt]
% [/#if]
\small
% [#assign sep = false]
% [#if (p.phone)?has_content]
{[=tex(p.phone)]}
% [#assign sep = true]
% [/#if]
% [#if (p.email)?has_content]
% [#if sep]
$|$
% [/#if]
\href{mailto:[=url(p.email)]}{\underline{{[=tex(p.email)]}}}
% [#assign sep = true]
% [/#if]
% [#if (p.linkedin)?has_content]
% [#if sep]
$|$
% [/#if]
\href{[=url(linkedinUrl(p.linkedin))]}{\underline{{[=tex(linkedinText(p.linkedin))]}}}
% [#assign sep = true]
% [/#if]
% [#if (p.location)?has_content]
% [#if sep]
$|$
% [/#if]
{[=tex(p.location)]}
% [/#if]
\end{center}

% [#if (p.summary)?has_content]
%-----------SUMMARY-----------
\noindent{\large \textsc{Summary}}\\[-4pt]
\begin{picture}(542,4)\put(0,2){\line(1,0){542}}\end{picture}\\
\hspace*{11pt}\parbox[t]{525pt}{\small {[=tex(p.summary)]}}

\vspace{6pt}

% [/#if]
% [#if education?has_content]
%-----------EDUCATION-----------
\noindent{\large \textsc{Education}}\\[-4pt]
\begin{picture}(542,4)\put(0,2){\line(1,0){542}}\end{picture}
% [#list education as edu]

\noindent\begin{picture}(536,13)
\put(11,4){\makebox(0,0)[l]{\bfseries {[=tex(edu.university!)]}}}
\put(536,4){\makebox(0,0)[r]{{[=tex(edu.location!)]}}}
\end{picture}
% [#if (edu.degree)?has_content || (edu.grade)?has_content || (edu.startYear)?has_content || (edu.endYear)?has_content]
\\
\begin{picture}(536,12)
\put(11,4){\makebox(0,0)[l]{\small\itshape {[=joinTex([edu.degree!, edu.grade!], "; ")]}}}
\put(536,4){\makebox(0,0)[r]{\small\itshape {[=joinTex([edu.startYear!, edu.endYear!], " -- ")]}}}
\end{picture}
% [/#if]
% [/#list]

\vspace{6pt}

% [/#if]
% [#if experience?has_content]
%-----------EXPERIENCE-----------
\noindent{\large \textsc{Experience}}\\[-4pt]
\begin{picture}(542,4)\put(0,2){\line(1,0){542}}\end{picture}
% [#list experience as exp]

\noindent\begin{picture}(536,13)
\put(11,4){\makebox(0,0)[l]{\bfseries {[=tex(exp.jobTitle!)]}}}
\put(536,4){\makebox(0,0)[r]{{[=joinTex([exp.startDate!, endLabel(exp)], " -- ")]}}}
\end{picture}
% [#if (exp.company)?has_content]
\\
\begin{picture}(536,12)
\put(11,4){\makebox(0,0)[l]{\small\itshape {[=tex(exp.company)]}}}
\end{picture}
% [/#if]
% [#assign items = bullets(exp.achievements!)]
% [#if items?has_content]
\begin{itemize}
\setlength{\itemsep}{0pt}\setlength{\parsep}{0pt}
\small
% [#list items as item]
\item {[=tex(item)]}
% [/#list]
\end{itemize}
% [#else]

\vspace{6pt}
% [/#if]
% [/#list]

% [/#if]
% [#if projects?has_content]
%-----------PROJECTS-----------
\noindent{\large \textsc{Projects}}\\[-4pt]
\begin{picture}(542,4)\put(0,2){\line(1,0){542}}\end{picture}
% [#list projects as proj]

\noindent\begin{picture}(536,13)
\put(11,4){\makebox(0,0)[l]{\small{\bfseries {[=tex(proj.name!)]}}
% [#if (proj.technologies)?has_content]
$|$ {\itshape {[=tex(proj.technologies)]}}
% [/#if]
}}
\end{picture}
% [#assign items = bullets(proj.description!)]
% [#if items?has_content]
\begin{itemize}
\setlength{\itemsep}{0pt}\setlength{\parsep}{0pt}
\small
% [#list items as item]
\item {[=tex(item)]}
% [/#list]
\end{itemize}
% [#else]

\vspace{6pt}
% [/#if]
% [/#list]

% [/#if]
% [#if certifications?has_content]
%-----------CERTIFICATES / ACHIEVEMENTS-----------
\noindent{\large \textsc{Certificates / Achievements}}\\[-4pt]
\begin{picture}(542,4)\put(0,2){\line(1,0){542}}\end{picture}
\begin{itemize}
\setlength{\itemsep}{0pt}\setlength{\parsep}{0pt}
\small
% [#list certifications as cert]
\item {{\bfseries {[=tex(cert.name!)]}}%
% [#if (cert.issuer)?has_content]
: {[=tex(cert.issuer)]}%
% [/#if]
}
% [/#list]
\end{itemize}

% [/#if]
% [#if skillRows?has_content]
%-----------TECHNICAL SKILLS-----------
\noindent{\large \textsc{Technical Skills}}\\[-4pt]
\begin{picture}(542,4)\put(0,2){\line(1,0){542}}\end{picture}\\
\small
% [#list skillRows as row]
\hspace*{11pt}{\bfseries {[=row.label]}}: {[=row.value]}
% [#sep]
\\
% [/#sep]
% [/#list]

% [/#if]
\end{document}
